\scanpage{193}

% B01 / text / top, continuation
\mspage{99}qui sont des structures de contour sur les
$E\in\unsure{\mathcal E|\mathscr E|\Sigma}$, \deleted{\unsure{ces}} structures \unsure{redonnant|redonnent} $A_i$ par

% B02 / display / upper, (13)
\[
 A_i=\bigcup_{E\in\mathcal E}A_{iE}\ . \tag*{(13)}
\]

% B03 / text / upper, before (14)
L'application $\mathrm{Et}_{A_i}$ correspondante

% B04 / diagram / upper, (14)
\diagram{(14)\qquad\begin{tikzcd}[ampersand replacement=\&,row sep=0pt]
 {\mathrm{Et}_{A_i}:}\& S\arrow[r,"\mathrm{surj}"]\&\mathcal E\\
 \&s\arrow[r,mapsto]\&\mathrm{Et}_{A_i}(s)
\end{tikzcd}}

% B05 / text / upper middle, following (14)
dépend a priori de la \deleted{\unsure{donnée|l'indice}} structure choisie $i\in\Pi$.

% B06 / text / middle, transitivity
\deleted{Notons que} Comme $G_0$ opère transitivement sur
$\mathcal E$ et respecte la structure $A_i$, \deleted{$A_i$}\inserted{les $A_{iE}$ \deleted{donc $A_i$}} sont connus,
donc $A_i$ est connu, dès qu'on connait
$A_{iE}$ pour \uline{un} $E_0\in\mathcal E$ :

% B07 / diagram / middle, (15)
\diagram{(15)\qquad L'application \inserted{$\unsure{\psi|\Psi|\varphi}_{E}:$}
\begin{tikzcd}[ampersand replacement=\&,row sep=small]
 \Pi\arrow[r]\&\substack{\text{Structures de contour}\\\text{sur }E_0}\\
 i\arrow[r,mapsto]\& A_{iE_0}=A_i\cap\unsure{P|\mathcal P|\mathfrak P}_2(E_0)
\end{tikzcd}}

% B08 / text / middle below (15)
est injective, et son image est formée
de structures de contour qui sont
\deleted{stables}\inserted{invariées} par l'opération de $G_{0E_0}$ (stabilisateur
de $E_0$ dans $G_0$) sur $E_0$.

% B09 / text / lower, system of contours
Il est clair d'ailleurs que la \deleted{\unsure{ens des}}\inserted{donnée d'une} structures
de contour sur $E_0$, invariée par $G_{E_0}$, est équivalente
à la donnée d'un système \inserted{$(C_E)_E$} de structures
de contour sur les $E\in\mathcal E$, invariées par $G$.

% B10 / text / lower, new paragraph
Il faut maintenant exprimer, en termes
\deleted{de la \unsure{s}} d'une

% B11 / display / foot, (16)
\[
 \begin{aligned}
 &\text{structure de contour }\unsure{C|{\mathbb C}}_{E_0}\text{ sur }E_0,\text{ \unsure{invarie|invariée}}\\
 &\text{par }G_{0E_0},
 \end{aligned} \tag*{(16)}
\]

% B12 / text / foot, last line
la condition qu'elle provient d'un $i\in\Pi$.

% archive sheet number: 192
